\documentclass[xcolor=table]{beamer}
% common packages

\usepackage[utf8]{inputenc}
\usepackage[T1]{fontenc}

\usepackage{aurical}
\usepackage{mathptmx}
\usepackage[scaled]{helvet}
\usepackage{lmodern}
\usepackage{textcomp}
\usepackage[scaled=.88]{DejaVuSansMono}

\usepackage[ngerman]{babel}
\usepackage[german=guillemets]{csquotes}

\usepackage{amsmath}
\usepackage{nicefrac}

\usepackage{tikz}
\usetikzlibrary{positioning,%
  shapes.geometric,%
  shapes.misc,%
  calc,%
  fit,%
  decorations,%
  decorations.pathreplacing,%
  backgrounds,%
  arrows,%
  automata,%
  trees%
}

\usepackage{pifont}
\newcommand{\goodmark}{\textcolor<presentation>{green!50!black}{\Pisymbol{pzd}{52}}}
\newcommand{\badmark}{\textcolor<presentation>{red}{\Pisymbol{pzd}{56}}}

% define colors used by presentation theme AND article
\xdefinecolor{bordeaux}{RGB}{128,0,50}
\xdefinecolor{niceblue}{RGB}{13,41,79}
\xdefinecolor{nicegreen}{RGB}{13,79,18}
\xdefinecolor{niceviolet}{RGB}{79,13,74}
\colorlet{maincolor}{orange}
\colorlet{alertedcolor}{red}
\colorlet{examplecolor}{green!50!black}
\usetheme{latex}

% color system
% - color!9 used as light fill color (e.g. for block content)
% - color!18 used as default fill color (e.g. for state)
% - color!30 used as highlight fill color (e.g. for block header)
% - color!50 used as strong highlight fill color (e.g. for header row)

\usepackage{listings}
\lstdefinelanguage[MyLaTeX]{TeX}[LaTeX]{TeX}%
  % TeX commands
  {moretexcs={enquote,includegraphics,%
    part,chapter,section,subsection,paragraph,subparagraph%
    tableofcontents,listoffigures,listoftables,maketitle,%
    subsection,subsubsection,paragraph,autoref,it,%
    textcolor,colorbox,xdefinecolor,colorlet,foreach,%
    rowcolors,rowcolor,lstdefinestyle,lstset,KOMAoptions,%
    setkomavar,setkomavar*,opening,closing,encl,%
    lehead,cehead,rehead,lefoot,cefoot,refoot,%
    lohead,cohead,rohead,lofoot,cofoot,rofoot,%
    ohead,chead,ihead,ofoot,cfoot,ifoot,%
    areaset,color,%
    automark,manualmark,markright,markboth,%
    tableofcontents,url,clearscrheadfoot,pagemark,headmark,%
    setheadtopline,setkomafont,setheadsepline,setfootsepline,%
    setfootbotline,chaptermark,thesection,thechapter,%
    thesubsection,subtitle,inst,section*,subsection*,institute},%
  % LaTeX environments
  morekeywords={[2]lstlisting,document,letter,center,flushleft,%
    flushright,align,itemize,enumerate,description,tabular,
    titlepage,figure,table,frame%
    },%
  % other things (like packages) to highlight
  morekeywords={[3]listings,textcomp,courier,xcolor,scrartcl,%
    scrlttr2,inputenc,babel,%
    fontenc,lmodern,mathptmx,%
    helvet,geometry,scrpage2,scrreprt,scrbook,%
    article,report,book,hyperref,%
    csquotes,amsmath,amssymb,la},
  alsoletter={0123456789*}
  }%

\lstdefinelanguage{BibTeX}
  {keywords={%
      @article,@book,@collectedbook,@conference,@electronic,@ieeetranbstctl,%
      @inbook,@incollectedbook,@incollection,@injournal,@inproceedings,%
      @manual,@mastersthesis,@misc,@patent,@periodical,@phdthesis,@preamble,%
      @proceedings,@standard,@string,@techreport,@unpublished%
      },
   comment=[l][\itshape]{@comment},
   morekeywords={[2]author,title,year,publisher,editor,booktitle,journal,series,volume,number,pages,institution,note,howpublished},
  }

\lstset{%
  basicstyle=\ttfamily,%
  language=[MyLaTeX]TeX,%
  texcsstyle=*\color{violet}\bfseries,%
  keywordstyle=\color{violet}\bfseries,%
  keywordstyle={[2]\color{green!70!black}\bfseries},%
  keywordstyle={[3]\color{maincolor!70!black}\bfseries},%
  commentstyle=\color{gray}\itshape,%
  stringstyle=\itshape,%
  numbers=none,%
  frame=lines,%
  backgroundcolor=\color{maincolor!10},%
  rulecolor=\color{maincolor!70},%
  framerule=1pt,%
  showstringspaces=false,%
  upquote=true,%
  framexleftmargin=3pt,%
  framexrightmargin=3pt}

\lstdefinestyle{pseudo}{language={},%
  basicstyle=\normalfont,%
  morecomment=[l]{//},%
  morekeywords={for,to,while,do,if,then,else},%
  mathescape=true,%
  columns=fullflexible}

\lstdefinestyle{block}{%
  backgroundcolor={},%
  frame=no,%
  aboveskip=0pt,%
  belowskip=0pt
}

% german umlauts and red braces
\lstset{
  literate={ö}{{\"o}}1
           {Ö}{{\"O}}1
           {ä}{{\"a}}1
           {Ä}{{\"A}}1
           {ü}{{\"u}}1
           {Ü}{{\"U}}1
           {ß}{{\ss}}1
           {\{}{{\textbf{\color{red}\{}}}1
           {\}}{{\textbf{\color{red}\}}}}1
           {[}{{\textbf{\color{red}[}}}1
           {]}{{\textbf{\color{red}]}}}1
}

\usepackage{tabularx}

% zebra tables
\newcommand{\mainrowcolors}{\rowcolors{1}{maincolor!30}{maincolor!9}}
\newenvironment{zebratabular}{\mainrowcolors\begin{tabular}}{\end{tabular}}
\newcommand{\setrownumber}[1]{\global\rownum#1\relax}
\newcommand{\headerrow}{\rowcolor{maincolor!50}\setrownumber1}

\AtBeginDocument{\hypersetup{
  pdftitle={\titlestring},
  pdfauthor={\authorstring}}
  \title[\shorttitlestring]{\titlestring}
  \author[\shortauthorstring]{\authorstring}
  \date{\datestring}}

% set texts used later
\newcommand{\targetsname}{Ziele dieses Vortrags}
\newcommand{\targetscontentname}{Ziele und Inhalt}
\newcommand{\outlinename}{Gliederung}

% often used sets
\newcommand{\set}[1]{\mathbb{#1}}
\newcommand{\R}{\set{R}}
\newcommand{\N}{\set{N}}
\newcommand{\Z}{\set{Z}}
\newcommand{\Q}{\set{Q}}

% easier use of operatorname
\newcommand{\op}[1]{\operatorname{#1}}

% display vectors in bold
\let\oldvec\vec
\let\vec\boldsymbol

\newcommand{\xxx}{\only<presentation>{\vskip1em}}

\newcommand{\inhead}[1]{\textbf{\color{maincolor}#1}}

\newenvironment{mybib}{%
  % temporary disable \section command to get rid of bibliography heading
  \renewcommand{\chapter}[2]{}
  \begin{thebibliography}{10}
}{%
  \end{thebibliography}
}

% load logo collection from "Die TeXnische Komödie"
\usepackage{dtklogos}

% load common user macros
\newcommand{\website}{
  \begin{Frame}[t]{Website}
    \alt<presentation>{%
      \vskip-4ex\hfill\parbox{6cm}{\centering
        \includegraphics[width=6cm]{images/qrcode}\\[1ex]
        \href{http://www.mlte.de/latex}{\huge\texttt{mlte.de/latex}}
      }
      \par\vskip2ex}{%
        Auf der Seite \url{http://www.mlte.de/latex} befinden sich}
    \begin{itemize}
      \item diese Präsentation, das Skript zum Vortrag,
      \item Beispieldokumente, Links zu weiteren Quellen und
      \item der Link zum Github-Repository\only<article>{.}
    \end{itemize}
  \end{Frame}
}

\newcommand*{\icon}[2][2]{%
  \tikz\draw[very thick, line join=round]
    (-.5,.5) -- ++(.7,0) -- ++(.3,-.3) --
    ++(-.3,0) -- ++(0,.3) -- ++(.3,-.3) -- ++(0,-.4)
    \foreach \x in {1,...,#1} {
      -- ++(0,-.2)
    }
    -- ++(-1,0) -- cycle
    (-.5,.2) node[fill, text=white,
      inner sep=2pt, minimum width=8mm, anchor=center] {#2}
    ++ (.1,-.3) -- ++(.8,0) \foreach \x in {1,...,#1} {
      -- ++(0,-.1) -- ++(-.8,0) -- ++(0,-.1) -- ++(.8,0)
    };
}

\colorlet{texicon}{examplecolor}
\colorlet{pdficon}{maincolor}
\colorlet{auxicon}{-examplecolor}
\colorlet{logicon}{black}
\colorlet{bblicon}{-maincolor}
\colorlet{bibicon}{alertedcolor}

\colorlet{tikzexample}{examplecolor!10}

\mode
<article>

\newcommand*{\examplebox}[2][center]{%
  \begingroup%
  \fboxsep=2ex%
  \begin{center}
    \colorbox{tikzexample}{#2}
  \end{center}
  \endgroup}

\mode
<presentation>

\newcommand*{\examplebox}[2][center]{%
  \begingroup%
  \fboxsep=2ex%
  \ifthenelse{\equal{#1}{center}}{%
    \hfill\colorbox{tikzexample}{#2}\hfill%
  }{%
    \colorbox{tikzexample}{#2}
  }%
  \endgroup}

\mode
<all>

\newcommand*{\tikzexample}[2][center]{\examplebox[#1]{%
  \begin{tikzpicture}#2\end{tikzpicture}}}

\newcommand*{\beamerexample}[2][1]{%
  \alt<presentation>{%
    \plain{\includegraphics[page=#1]{#2}}
  }{
    \plain{\frame{\includegraphics[width=10cm,page=#1]{#2}}}
  }
}

\newcommand{\currentauthor}{Malte}
\newcommand{\updateauthorcolors}{\ifthenelse{\equal{\currentauthor}{Malte}}{
  \colorlet{authormalte}{maincolor}\colorlet{authorjonny}{black}
}{
  \colorlet{authormalte}{black}\colorlet{authorjonny}{maincolor}
}}
\newcommand{\insertcurrentauthor}{\currentauthor}
\updateauthorcolors
\newcommand{\malte}{\renewcommand{\currentauthor}{Malte}\updateauthorcolors}
\newcommand{\jonny}{\renewcommand{\currentauthor}{Johannes}\updateauthorcolors}

\newcommand{\lorem}{\textcolor{black!40}{Auch in der zweiten Zeile. Lorem ipsum dolor sit amet, consectetur, adipisci velit, \ldots}}

\newcommand{\colorsample}[1]{\textcolor{#1}{\rule[-.5ex]{2em}{2ex}} \texttt{#1}}








% set options for article
\mode
<article>

\newcommand{\itemref}[1]{#1}
\newcommand{\term}[1]{#1}

\KOMAoption{fontsize}{12pt}
\KOMAoption{parskip}{half}
\KOMAoption{open}{any}
\KOMAoption{headings}{big}
\KOMAoption{numbers}{noendperiod}
\KOMAoption{twoside}{false}
\KOMAoption{mpinclude}{true}

\usepackage{paralist}
% use compactitem as default itemize environment
\newenvironment{looseitemize}{}{}
\let\looseitemize\itemize
\let\endlooseitemize\enditemize
\let\itemize\compactitem
\let\enditemize\endcompactitem
% use compactenum as default enumerate environment
\newenvironment{looseenumerate}{}{}
\let\looseenumerate\enumerate
\let\endlooseenumerate\endenumerate
\let\enumerate\compactenum
\let\endenumerate\endcompactenum

% create head and foot line
\usepackage{scrpage2}
\pagestyle{scrheadings}

% stretch headline above margin (scrguide -- page 244)
\setheadwidth[0pt]{textwithmarginpar}

% create targets block
\usepackage{minitoc}
\nomtcrule
\ktightmtctrue
\tightmtctrue
\setlength{\mtcindent}{0pt}
\mtcsetfont{minitoc}{*}{\normalsize\rmfamily\upshape\mdseries}
\mtcsetfont{minitoc}{section}{\normalsize\rmfamily\upshape\mdseries}
\mtcsetfont{minitoc}{subsection}{\normalsize\rmfamily\upshape\mdseries}
\mtcsettitlefont{minitoc}{\normalsize\sffamily\upshape\bfseries}

\makeatletter
\newcommand{\targets}[1]{
  \parbox[t]{0.46\textwidth}{\raggedright
    \kern-0.8\baselineskip\nopagebreak[4]%
    \par\noindent %%
    \begin{tabular}{@{}p{\columnwidth}@{}}
    \reset@font\mtifont\do@mtitc{\mtc@v\targetsname}\\
    \end{tabular}%
    \begin{mtc@verse}{\mtcoffset}
    \begin{enumerate}#1\end{enumerate}
    \end{mtc@verse}
  }
  \hskip0.04\textwidth
  \parbox[t]{0.48\textwidth}{\raggedright
    \minitoc
  }
  % emulate two frames
  \addtocounter{slidenumber}{2}
  \addtocounter{framenumber}{2}
}
\makeatother

% keep the name of the current chapter and the current section

\clearscrheadfoot % clear all default settings
\automark[chapter]{section} % \automark[\rightmark]{\leftmark}

\ihead{\shorttitlestring\ \rightmark\ \textcolor{maincolor}{\bfseries |} \leftmark \hfill \pagemark}
\setkomafont{pageheadfoot}{\normalfont\sffamily}
\setkomafont{pagenumber}{\usekomafont{pageheadfoot}\color{maincolor}}
\setheadsepline{0.5pt}[\color{maincolor}]


% Create own counter as replacement for framenumber. framenumber is
% incremented in article mode as \refstepcounter{framenumber} by
% beamer and hyperref complains about duplicated targets if framenumber
% gets resetted after each chapter.
\newcounter{slidenumber}

\makeatletter
% create AtBeginChapter analog to AtBeginDocument
\def\AtBeginChapter{\g@addto@macro\@beginchapterhook}

\@onlypreamble\AtBeginChapter

\let\oldchapter\chapter
\long\def\chapter{\@beginchapterhook
  \oldchapter}
  
\ifx\@beginchapterhook\@undefined
  \let\@beginchapterhook\@empty
\fi
\makeatother

% slidenumbers starts over with every new chapter with 2 on chapterframe
\AtBeginChapter{\setcounter{slidenumber}{2}}

% hide block title in article mode
\setbeamertemplate{block begin}{}
\setbeamertemplate{block end}{}

% display block title in Block environment
\newenvironment{Block}{%
  \setbeamertemplate{block begin}[default]
  \setbeamertemplate{block end}[default]
  \begin{block}}{\end{block}}

% hide frame title and subtitle in article mode for default frame environment
\setbeamertemplate{frametitle}{}

% display frame title of Frame environment
\newenvironment{Frame}{%
  \setbeamertemplate{frametitle}{%
    \paragraph*{%
      \insertframetitle%
      \ifx\insertframesubtitle\empty\else%
        \newline%
        \mdseries\insertframesubtitle%
      \fi%
    }\ \par}
  \begin{frame}}{\end{frame}}

% hide effect of beameritemize
\newenvironment{beameritemize}{%
  \let\item\relax\ignorespaces}{%
  \ignorespacesafterend}
  
% map beamerenumerate to inparaenum
\newenvironment{beamerenumerate}{%
  \begin{inparaenum}\ignorespaces}{%
  \end{inparaenum}\ignorespacesafterend}

% map plain command to frame
\newcommand{\plain}[1]{\frame{#1}}

% display frame number
%\setbeamertemplate{frame end}{%
%  \marginpar{\hfill\scriptsize\sffamily\color{maincolor}$\to\,$\insertframenumber}}

% display framenumber as chapter number and slide number
\renewcommand{\insertframenumber}{\thechapter-\theslidenumber}

% incremend slide number
\setbeamertemplate{frame begin}{\stepcounter{slidenumber}}
% display slide number
\setbeamertemplate{frame end}{%
  \marginpar%
    [\scriptsize\sffamily\color{maincolor}\insertframenumber$\,\leftarrow$\hfill]%
    {\hfill\scriptsize\sffamily\color{maincolor}$\rightarrow\,$\insertframenumber}}
      
\newcommand{\script}{\addtomargin{Script}\mode<article>}
      
\usepackage{hyperref}
\hypersetup{breaklinks=true,
            pdfborder={0 0 0},
            pdfhighlight={/N}}

\usepackage{verbatim} % for the comment environment
















% Set options for beamer
\mode
<presentation>

\makeatletter
% Reverts "Heiko's fix for correct generation of section* and
% subsection* bookmarks" in beamerbasecompatibility.sty which
% disables bookmarks in \section* and \subsection*.
% This is done by disabling \Hy@writebookmark in \beamer@section
% and \beamer@subsection if their second argument is empty.
% This is the case for \section* and \subsection* as these
% commands are mapped to \beamer@section[{#1}]{} resp.
% \beamer@subsection[{#1}]{}. By defining \beamer@section and
% \beamer@subsection again to the original and not patched version
% in the macro \beamer@lastminutepatches where the patch is applied
% this behaviour gets disabled. This is mainly reverted because
% bookmarks of parts are not added after \section*{Conclusion} due
% to this patch.
\g@addto@macro\beamer@lastminutepatches{
  \ifx\org@beamer@section\undefined{}%
  \else
    \let\beamer@section\org@beamer@section{}%
  \fi
  \ifx\org@beamer@subsection\undefined{}%
  \else
    \let\beamer@subsection\org@beamer@subsection{}%
  \fi
}
\makeatother

% display framenumber as part number and frame number
\renewcommand{\insertframenumber}{\thepart-\theframenumber}

% reset frame number at the beginning of each chapter
\AtBeginPart{
  \setcounter{framenumber}{0}
}

\AtBeginSection[]{
  \begingroup
  \setbeamercolor{background canvas}{bg=maincolor!9}
  \plain{
    \Huge\color{maincolor}\bfseries
    \vskip11ex
    \begin{center}
      \insertsection
    \end{center}
    \vfill
    \normalsize
    \hfill\insertcurrentauthor\ \strut
  }
\endgroup
}

\newcommand{\chapter}[1]{
  \part{#1}
  \title{\LaTeX\\#1}
  \plain{
  \begin{center}
    \vskip4ex
    \textbf{\fontsize{36pt}{36pt}\selectfont\rmfamily\color{maincolor}\LaTeX}
    \par\vskip2ex
    \large Johannes und Malte auf der MetaNook 2013
  \end{center}

  \vskip2ex\hskip4em
  \begin{minipage}{20em}\Large
    18:00 Uhr\\
    \tikz[remember picture] \node[coordinate] (n1) {};
    \ \textbf{\color{maincolor}Grundlagen}

    \vskip1ex
    19:15 Uhr\\
    \tikz[remember picture] \node[coordinate] (n2) {};
    \ \textbf{\color{maincolor}Fortgeschrittene Verwendung}

    \vskip1ex
    20:45 Uhr\\
    \tikz[remember picture] \node[coordinate] (n3) {};
    \ \textbf{\color{maincolor}Präsentieren mit {\normalsize BEAMER}}

    \vskip1ex
    21:45 Uhr\\
    \tikz[remember picture] \node[coordinate] (n4) {};
    \ \textbf{\color{maincolor}Ti\emph{k}Z}
  \end{minipage}

  \begin{tikzpicture}[remember picture,overlay]
    \draw[line width=7pt,maincolor,->]
      (n\insertpartnumber) ++(-6em,1ex) -- ++(6em,0);
  \end{tikzpicture}
}

}

\newcommand{\itemref}[1]{\textbf{\color{maincolor}#1}}

\newcommand{\term}[1]{{\color{blue}#1}}
  
\newcommand{\targets}[1]{\section*{\targetscontentname}
  \frame[plain]{\frametitle{\targetsname}%
  \begin{enumerate}#1\end{enumerate}}
  \frame[plain]{\frametitle{\outlinename}\tableofcontents}}

% display math in serif font on the slides
\usefonttheme[onlymath]{serif}

% create new environment Block behaving exactly like block
\newenvironment{Block}{%
  \begin{block}}{\end{block}}

% create new environment Frame behaving exactly like frame in beamer mode
\newenvironment{Frame}[1][]{%
  \begin{frame}[environment=Frame,#1]}{\end{frame}}

% create new environment beameritemize behaving exactly like itemize
\newenvironment{beameritemize}{%
  \begin{itemize}}{\end{itemize}}

% create new environment beamerenumerate behaving exactly like enumerate
\newenvironment{beamerenumerate}{%
  \begin{enumerate}}{\end{enumerate}}


\mode
<all>


\usepackage{multicol}

\newcommand{\authorstring}{Johannes und Malte}
\newcommand{\shortauthorstring}{\textcolor{authormalte}{Malte} \& \textcolor{authorjonny}{Johannes}}
% author colors are updated in macros.tex

\newcommand{\titlestring}{\LaTeX}
\newcommand{\shorttitlestring}{\LaTeX}
\newcommand{\datestring}{MetaNoOK 2015}

\begin{document}
\jonny

\label{chapter-beamer}
\setcounter{part}{2}
\chapter{Präsentationen mit \beamer}

\targets{
  \item \strut\beamer\ verwenden können.
  \item Vor- und Nachteile von \beamer\ kennen und einschätzen können, wann und wofür \beamer\ gut geeignet ist.
  \item Fortgeschrittene Anwendungsmöglichkeiten von \beamer\ kennen lernen.
}

\website

\beamersection{Was ist \beamer?}

\subsection{Einleitung}

\begin{Frame}{Was ist \beamer?}
  \begin{itemize}
    \item \alert{Dokumentenklasse für \LaTeX} für die Erzeugung von Präsentationen.
      \only<article>{\newline(Diese Präsentation und das Skript wurden mit \beamer\ erzeugt.)}
    \item Keine eigene und \alert{keine graphische Anwendung}.
    \item \strut\beamer\ ist in \MiKTeX\ und \TeX\ Live enthalten.\newline
      (\alert{Es kann direkt losgehen}.)
  \end{itemize}
\end{Frame}

\mode
<article>

\begin{Frame}{Historie}
  \begin{description}
    \item[1998] Till Tantau erzeugt sich erste Makros für Präsentationen.
    \item[2003] Er verwendet die erste Version für seine Promotionsverteidigung.
    \item[2003] Veröffentlichung und Implementierung von vielen Benutzerwünschen.
    \item[2007] \beamer\ wird nicht weiter gepflegt.
    \item[2010] \beamer\ wird an Joseph Wright and Vedran Mileti\'c übergeben.
    \item[2014] Aktuelle Version 3.33 wird kontinuierlich weiter entwickelt.
  \end{description}
\end{Frame}

\mode
<all>

\begin{Frame}{Workflow}
  \lstset{language={}}
  \begin{enumerate}
    \item Normales \LaTeX-Dokument erzeugen.\\
      Dabei einige spezielle \beamer-Kommandos verwenden.
    \item \LaTeX-Dokument mit \lstinline-pdflatex- kompilieren.
    \item Ergebnis überprüfen und \LaTeX-Dokument anpassen.
  \end{enumerate}
\end{Frame}

\subsection{Eigenschaften}

\begin{Frame}{Funktionsweise von \beamer}
  \begin{itemize}
    \item Kompilieren wie jedes andere \LaTeX-Dokument auch.
    \item Normale \LaTeX-Kommandos funktionieren.
    \item Sinnvolles funktionales Aussehen von Vorträgen.
    \item Einfaches Ein- und Ausblenden von Seitenteilen.
    \item Automatische Gliederungen und Navigationsleisten.
    \item Präsentationen im PDF-Format können auf jedem Computer dargestellt werden.
  \end{itemize}
\end{Frame}

\begin{Frame}{\beamer\ vs. PowerPoint}
  \begin{zebratabular}{rcc}
    \headerrow Aspekte & \beamer\ & PowerPoint \\
    Erlernen ohne \LaTeX-Kenntnisse & \badmark\badmark & \goodmark \\
    Objekte frei positionieren & \badmark & \goodmark\goodmark \\
    Grafiken direkt erstellen & \badmark & \goodmark \\
    Einbinden von Multimedia & -- & \goodmark \\
    Arbeitsgeschwindigkeit Anfänger & -- & -- \\
    Arbeitsgeschwindigkeit Profi & \goodmark & \goodmark \\
    Erlernen mit \LaTeX-Kenntnissen & \goodmark & \goodmark \\
    Dokumentation & \goodmark & \goodmark \\
    Vorlagenqualität & \goodmark & -- \\
    Typographie & \goodmark & \badmark\badmark \\
    Konsistenz des Aussehens & \goodmark\goodmark & \badmark \\
    Visualisierung des Vortragsaufbaus & \goodmark\goodmark & \badmark \\
    Mathematische Formeln & \goodmark\goodmark & \badmark\badmark \\
    Quelltextdarstellung & \goodmark\goodmark & \badmark\badmark
  \end{zebratabular}
\end{Frame}


\beamersection{Verwendung von \beamer}

\begin{Frame}[fragile]{Beispiel}
  \begin{lstlisting}[gobble=4]
    \documentclass{beamer}

    \usepackage[utf8]{inputenc}
    \usepackage[T1]{fontenc}
    \usepackage{lmodern}
    \usepackage[ngerman]{babel}

    \begin{document}
      \begin{frame}{Funktionen von Beamer}
        Kompilieren wie jedes andere
        \LaTeX-Dokument auch.
      \end{frame}
    \end{document}
  \end{lstlisting}
\end{Frame}

\beamerexample{demo/beamer.pdf}

\subsection{Folien}

\begin{Frame}[fragile]{Folien}
  \begin{itemize}
    \item Ein \beamer-Dokument besteht aus Folien.
    \item Die Umgebung \lstinline-frame- verarbeitet
      bis zu zwei Parameter in gescheiften Klammern \lstinline-{}-
    \item Der erste Parameter ist der Folientitel.
    \item Der zweite Parameter ist der Untertitel.
    \item Innerhalb der Umgebung \lstinline|frame| wird normaler \LaTeX-Code
      verwendet.
  \end{itemize}
\end{Frame}

\begin{Frame}[fragile]{Titelfolie}
  \begin{Block}{In der Präambel}
    \begin{lstlisting}[gobble=6,style=block]
      \title[Kurztitel]{%
        Lange Version des langen Titels}
      \subtitle{Ein langer Untertitel beschreibt
        alles noch etwas genauer.}
      \author[Thorn, Schmitz]{%
        Johannes Thorn \and Malte Schmitz}
      \date[KPT 2014]{Konferenz über
        Präsentationstechniken, 2014}
    \end{lstlisting}
  \end{Block}

  \begin{lstlisting}[gobble=4]
    \begin{frame}
      \maketitle
    \end{frame}
  \end{lstlisting}
\end{Frame}

\begin{Frame}[fragile]{Angabe von Instituten}
  \begin{lstlisting}[gobble=4]
    \author[Thorn, Schmitz]{%
      Johannes Thorn\inst{1}
      \and Malte Schmitz\inst{2}}

    \institute[Hier und Dort]{%
      \inst{1}Ein Institut\\
      Universität Hier
      \and
      \inst{2}Noch ein Institut\\
      Universität Dort}
  \end{lstlisting}
\end{Frame}

\beamerexample{demo/beamer-titlepage.pdf}

\begin{Frame}[fragile]{Inhaltsverzeichnis}
  \begin{itemize}
    \item Strukturbefehle außerhalb von \lstinline-frame-\\
      normal verwenden.
      \begin{itemize}
        \item ca. 3 Abschnitte mit \lstinline-\section-
        \item je max. 4 Unterabschnitte mit \lstinline-\subsection-
      \end{itemize}
    \item \lstinline-\tableofcontents- im \lstinline-frame- setzt das Inhaltsverzeichnis.
    \item Je nach Theme erscheinen \lstinline-\section- und
      \lstinline-\subsection- auch in Navigationsleisten.
    \item \lstinline-\section*- und \lstinline-\subsection*- erscheinen in
      Navigationsleisten aber nicht im Inhaltsverzeichnis.
  \end{itemize}
\end{Frame}

\subsection{Strukturelemente}

\begin{Frame}[fragile]{Listen, Tabellen und Grafiken}
  \begin{itemize}
    \item Listen mit \lstinline-itemize- und \lstinline-enumerate-,
    \item Tabellen mit \lstinline-tabular- und
    \item Grafiken mit \lstinline-\includegraphics- funktionieren wie immer in \LaTeX.
  \end{itemize}

  \xxx

  \begin{itemize}
    \item Eine Folie ist 128~mm $\times$ 96~mm groß.
    \item Zeilenumbruch \lstinline-\\- zum Ausrichten von Text sinnvoll.
  \end{itemize}
\end{Frame}

\begin{Frame}[fragile]{Formelsatz}
  \begin{itemize}
    \item Formelsatz wie immer in \LaTeX
    \item zum Beispiel \lstinline-$1+1=2$-\\
      oder \lstinline-\[1+1=2\]-
  \end{itemize}

  \xxx

  \begin{lstlisting}[gobble=4]
    % Formeln mit Serifen setzen
    \usefonttheme[onlymath]{serif}
  \end{lstlisting}
\end{Frame}

\begin{Frame}[fragile]{Blöcke}
  \begin{lstlisting}[gobble=4]
    \begin{block}{Überschrift}
      Dieser Text steht im normalen Block.
    \end{block}

    \begin{alertblock}{Achtung}
      Dieser Text steht im hervorgehobenen Block.
    \end{alertblock}

    \begin{exampleblock}{Beispiel}
      Dieser Text steht im Beispielblock.
    \end{exampleblock}
  \end{lstlisting}
\end{Frame}

\beamerexample{demo/beamer-struktur.pdf}

\begin{Frame}[fragile]{Theorem-Umgebungen}
  \begin{lstlisting}[gobble=4]
    \begin{Satz}[Sandhaufensatz]
      Es gibt keine Sandhaufen.
    \end{Satz}

    \begin{Beweis}
      \begin{enumerate}
        \item Ein Sandkorn ist kein Sandhaufen.
        \item Sandkörner werden durch Hinzufügen
          eines Sandkorns nicht zum Sandhaufen.
        \item Induktiv folgt die Aussage. \qedhere
      \end{enumerate}
    \end{Beweis}

    \begin{Beispiel}
      Vergleiche unsere Baustellen.
    \end{Beispiel}
  \end{lstlisting}
\end{Frame}

\beamerexample[2]{demo/beamer-struktur.pdf}

\begin{Frame}[fragile]{Spalten}
  \begin{lstlisting}[gobble=4]
    \begin{frame}{Spalten}
      \begin{columns}
        \begin{column}{5cm}
          Linke Spalte.
        \end{column}
        \begin{column}{5cm}
          Rechte Spalte.
        \end{column}
      \end{columns}
    \end{frame}
  \end{lstlisting}
\end{Frame}

\beamerexample[3]{demo/beamer-struktur.pdf}

\begin{Frame}[fragile]{Quelltext ist fragil.}
  \begin{Block}{In der Präambel}
    \begin{lstlisting}[style=block,gobble=6]
      \usepackage{listings}
      \lstset{%
        basicstyle=\ttfamily,%
        showstringspaces=false,%
        upquote=true}
      \usepackage{textcomp} % für upquote
      \usepackage{courier} % für schönere Schriftart
    \end{lstlisting}
  \end{Block}

  \xxx

  \lstinputlisting{demo/beamer-listing-content.tex}
\end{Frame}

\beamerexample{demo/beamer-listing.pdf}

\subsection{Form}

\begin{Frame}[fragile]{Themes}
  \begin{Block}{Theme}
    Wird geladen mit \lstinline-\usetheme{name}- und bestimmt die
    \alert{allgemeine Form} der Präsentation.
  \end{Block}

  \begin{Block}{Inner Theme}
    Wird geladen mit \lstinline-\useinnertheme{name}- und bestimmt
    die \alert{Form des Folieninhalts}.
  \end{Block}

  \begin{Block}{Outer Theme}
    Wird geladen mit \lstinline-\useoutertheme{name}- und bestimmt
    die \alert{Form der Layoutelemente}.
  \end{Block}

  \begin{Block}{Color Theme}
    Wird geladen mit \lstinline-\usecolortheme{name}- und bestimmt
    die \alert{allgemeine Farbe} der Präsentation.
  \end{Block}
\end{Frame}

\beamerexample[23]{demo/beamer-Boadilla.pdf}

\beamerexample[23]{demo/beamer-Madrid.pdf}

\beamerexample[23]{demo/beamer-Rochester.pdf}

\beamerexample[23]{demo/beamer-Montpellier.pdf}

\beamerexample[23]{demo/beamer-Goettingen.pdf}

\beamerexample[23]{demo/beamer-Frankfurt.pdf}

\beamerexample[23]{demo/beamer-Luebeck.pdf}

\begin{Frame}{Themes Matrix}
  \begin{itemize}
    \item Das war nur eine kleine Auswahl der möglichen Kombinationen.
    \item Die vollen Variationsmöglichkeiten ergeben sich erst aus der
      Kombination von Theme, Inner Theme, Outer Theme und Color Theme.
  \end{itemize}

  \xxx

  \begin{mybib}
    \bibitem{themes}
      Sebastian Pipping.
      \newblock The \beamer\ Theme Matrix.
      \newblock \alt<presentation>{\href{http://www.hartwork.org/beamer-theme-matrix/}{\texttt{hartwork.org/beamer-theme-matrix}}}{\url{http://www.hartwork.org/beamer-theme-matrix/}}, April 2009.
  \end{mybib}
\end{Frame}

\mode
<article>

\begin{Block}{Navigationssymbole ausblenden}
  \begin{lstlisting}[gobble=4]
    % hide navigation symbols
    \setbeamertemplate{navigation symbols}{}
  \end{lstlisting}
\end{Block}

\mode
<all>


\beamersection{Fortgeschrittene Verwendung}

\subsection{Overlays}

\begin{Frame}[fragile]{Einfache Overlays}
  Kommando \lstinline-\pause- blendet Elemente schrittweise ein.

  \begin{lstlisting}[gobble=4]
    \begin{enumerate}
      \item Ein Sandkorn ist kein Sandhaufen.
        \pause
      \item Sandkörner werden durch Hinzufügen
        eines Sandkorns nicht zum Sandhaufen.
        \pause
      \item Induktiv folgt die Aussage.
    \end{enumerate}
  \end{lstlisting}

  \xxx

  \begin{onlyenv}<presentation>
    \begin{enumerate}
      \item Ein Sandkorn ist kein Sandhaufen.
        \pause
      \item Sandkörner werden durch Hinzufügen
        eines Sandkorns nicht zum Sandhaufen.
        \pause
      \item Induktiv folgt die Aussage.
    \end{enumerate}
  \end{onlyenv}
\end{Frame}

\begin{Frame}[fragile]{Overlay-Spezifikationen}
  \begin{lstlisting}[gobble=4]
    \begin{Satz}[Sandhaufensatz]
      Es gibt keine Sandhaufen.
    \end{Satz}

    \begin{Beweis}<2->
      \begin{enumerate}
        \item<3-> Ein Sandkorn ist kein Sandhaufen.
        \item<4-> Sandkörner werden durch Hinzufügen
          eines Sandkorns nicht zum Sandhaufen.
        \item Induktiv folgt die Aussage. \qedhere
      \end{enumerate}
    \end{Beweis}

    \onslide<5->

    Der \alert<6>{Induktionsbeweis} ist
    \alert<7>{falsch}!
  \end{lstlisting}
\end{Frame}

\mode
<presentation>

\begin{Frame}{Overlay-Spezifikationen}
  \begin{Satz}[Sandhaufensatz]
    Es gibt keine Sandhaufen.
  \end{Satz}

  \begin{Beweis}<2->
    \begin{enumerate}
      \item<3-> Ein Sandkorn ist kein Sandhaufen.
      \item<4-> Sandkörner werden durch Hinzufügen
        eines Sandkorns nicht zum Sandhaufen.
      \item Induktiv folgt die Aussage. \qedhere
    \end{enumerate}
  \end{Beweis}

  \onslide<5->

  Der \alert<6>{Induktionsbeweis} ist
  \alert<7>{falsch}!
\end{Frame}

\mode
<article>

\frame{}

\mode
<all>

\begin{Frame}[fragile]{Ein- und Ausblenden}
  \begin{itemize}
    \item \lstinline|\uncover<3->{Inhalt}| blendet Inhalt erst ab
      Folie~3 ein. Der Platz wird jedoch vorher schon belegt.
    \item \lstinline|\only<3->{Inhalt}| setzt Inhalt erst ab Folie~3.
      Zuvor wird kein Platz belegt.
  \end{itemize}

  \xxx
  \pause

  \begin{lstlisting}[gobble=4]
    In diesem \uncover<2->{Satz} werden
    \only<3->{Worte }eingeblendet.
  \end{lstlisting}
  \only<presentation>{In diesem \uncover<3->{Satz} werden
    \only<4->{Worte }eingeblendet.}
\end{Frame}

\subsection{Artikelfassung}

\begin{Frame}[fragile]{Artikelfassung}
  \begin{Block}{Ziel}
    Generierung von Artikelfassung und Präsentation
    aus demselben Quellen-Dokument.
  \end{Block}

  \xxx
  \pause

  \begin{alertblock}{Problem}
    \begin{tabular}{r@{ }l}
      \textbf{Folien} & Dokumentenklasse von \beamer.\\
      \textbf{Artikel} & Dokumentenklasse von KOMA-Script.
    \end{tabular}
  \end{alertblock}

  \xxx
  \pause

  \begin{Block}{Lösung}
    \begin{itemize}
      \item Zwei \LaTeX-Dokumente für beide Dokumentenklassen.
      \item Drittes \LaTeX-Dokument für den Inhalt.
      \item Einbinden des Inhalts mit \lstinline-\input-.
    \end{itemize}
  \end{Block}
\end{Frame}

\begin{Frame}[fragile]{Einbinden des Inhalts}
  \begin{center}
    \begin{tikzpicture}[
        on grid,
        auto,
        node distance=18mm and 30mm,
        engine/.style={
          font=\rmfamily\Large\bfseries,
          inner sep=2pt
        }
      ]

      \node (content tex) {\shortstack{\textcolor{texicon}{\icon{TEX}}\\\texttt{content.tex}}};

      \uncover<2->{
        \node[below left=of content tex] (beamer tex)
          {\shortstack{\textcolor{texicon}{\icon{TEX}}\\\texttt{beamer.tex}}};
        \node[below right=of content tex] (article tex)
          {\shortstack{\textcolor{texicon}{\icon{TEX}}\\\texttt{article.tex}}};
      }

      \uncover<3->{
        \node[below=of beamer tex, engine]
          (make beamer) {\LaTeX.mk};

        \node[below=of article tex, engine]
          (make article) {\LaTeX.mk};

        \node[below=of make beamer] (beamer pdf)
          {\shortstack{\textcolor{pdficon}{\icon{PDF}}\\\texttt{beamer.pdf}}};
        \node[below=of make article] (article pdf)
          {\shortstack{\textcolor{pdficon}{\icon{PDF}}\\\texttt{article.pdf}}};
      }

      \only<2->{
        \draw[very thick]
        (content tex.west) edge[->] node[swap, near start]
          {\texttt{\color{texcs}\bfseries\textbackslash input}} (beamer tex.north);

        \draw[very thick]
          (content tex.east) edge[->] node[near start]
            {\texttt{\color{texcs}\bfseries\textbackslash input}} (article tex.north);
      }

      \only<3->{
        \draw[very thick]
          (beamer tex) edge[->] (make beamer)
          (article tex) edge[->] (make article);

        \draw[very thick]
          (make beamer) edge[->] (beamer pdf)
          (make article) edge[->] (article pdf);
      }
    \end{tikzpicture}
  \end{center}
\end{Frame}

\begin{Frame}[fragile]{Inhalt \texttt{content.tex}}
  \begin{lstlisting}[gobble=4]
    \usepackage[utf8]{inputenc}
    \usepackage[T1]{fontenc}
    \usepackage{lmodern}
    \usepackage[ngerman]{babel}

    \title{Mein Vortrag}
    \author{Mein Name}

    \begin{document}
      \begin{frame}
        \maketitle
      \end{frame}

      \begin{frame}{Folientitel}
        Hier passierts \dots
      \end{frame}
    \end{document}
  \end{lstlisting}
\end{Frame}

\begin{Frame}[fragile]{Dokumentenklassen}
  \inhead{Für die Folien \texttt{beamer.tex}}
  \begin{lstlisting}[gobble=4]
    % Beamer als Dokumentenklasse verwenden
    \documentclass{beamer}
    % gemeinsamen Inhalt einbinden
    \input{content.tex}
  \end{lstlisting}

  \xxx

  \inhead{Für den Artikel \texttt{article.tex}}
  \begin{lstlisting}[gobble=4]
    % KOMA-Script als Dokumentenklasse verwenden
    \documentclass{scrartcl}
    % Einstellungen für KOMA-Script
    \KOMAoptions{parskip=full}
    % Beamer als Paket laden
    \usepackage{beamerarticle}
    % gemeinsamen Inhalt einbinden
    \input{content.tex}
  \end{lstlisting}
\end{Frame}

\begin{Frame}[fragile]{Modes}
  \begin{tabular}{r@{ }l}
    \texttt{presentation} & nur für Folien\\
    \texttt{article} & nur für Artikel\\
    \texttt{all} & für Folien und Artikel (Standard)
  \end{tabular}

  \vskip3ex

  \begin{lstlisting}[gobble=4]
    \mode
    <name>
  \end{lstlisting}
  Wechselt den aktuellen Mode.

  \xxx

  \begin{lstlisting}[gobble=4]
    \mode*
  \end{lstlisting}
  Automatische Modeumschaltung:
  \begin{itemize}
    \item Innerhalb von \lstinline-frame- Mode {all}.
    \item Außerhalb von \lstinline-frame- Mode {article}.
  \end{itemize}
\end{Frame}


\section*{Zusammenfassung}

\begin{frame}[fragile]{Zusammenfassung}
  \begin{enumerate}
    \item Mit der Dokumentenklasse \lstinline-beamer- können \alert{sehr
          leicht Präsentationen erstellt} werden, wenn man mit \LaTeX\ etwas geübt ist.
    \item Folien werden mit der Umgebung \lstinline-frame- erzeugt.
      Fast alle \alert{\LaTeX-Kommandos funktionieren wie immer}.
    \item Mit \alert{Listen, Blöcken, Theoremen und Spalten} wird
      der Inhalt auf den Folien \alert{strukturiert}.
    \item \alert{Overlay- und Mode-Spezifikationen} werden in spitzen
      Klammern \lstinline-<- und \lstinline->- angegeben. Diese beeinflussen, in welchem
      \alert{Schritt der Animation} und in welchem \alert{Mode}
      das Kommando ausgeführt wird.
    \item Ein \LaTeX-Dokument, das Folien enthält, kann auch \alert{als Artikel kompiliert} werden.
    \item Bei Problemen \alert{lies die Anleitung.}
  \end{enumerate}
\end{frame}

\begin{Frame}{Zum Weiterlesen}
  \begin{mybib}
    \bibitem{Tantau}
      Till Tantau, Joseph Wright und Vedran Mileti\'c.
      \newblock The \textsc{beamer} \textit{class}, User Guide.
      \newblock \alt<presentation>{\href{http://mirrors.ctan.org/macros/latex/contrib/beamer/doc/beameruserguide.pdf}{\texttt{beameruserguide.pdf}}}{\url{http://mirrors.ctan.org/macros/latex/contrib/beamer/doc/beameruserguide.pdf}}, Oktober 2013.

    \bibitem{Tantau}
      Till Tantau.
      \newblock \emph{Beamer: Strahlende Vorträge mit \LaTeX},
      \newblock Präsentieren und Dokumentieren -- Tools.
      \newblock Vorlesung vom 31. Oktober 2012.
  \end{mybib}
\end{Frame}


\website

\end{document}

